%----------------------------------------------------------------------
% 實驗設計分析
%----------------------------------------------------------------------

\chapter{Experimental Results}\label{chap:experimental_results}

\section{Experimental Setup and Evaluation Metrics}\label{sec:4-setup}

The main experiments were conducted using a 1.3B-parameter MTP model
with four prediction heads. The models were trained on
OpenWebText\cite{gokaslan2019openwebtext}. The baseline and proposed
models were trained on 10 billion tokens, while the ablation experiments
used models trained on 1 billion tokens. The main SDC+LCM configuration
used $\lambda_{\mathrm{SDC}}=1.0$, $\tau=0.3$, and
$\lambda_{\mathrm{LCM}}=0.5$. Detailed training configurations are
provided in Appendix \ref{app:experimental-config}.

The baseline model was optimized using the original MTP training objective,
whereas the proposed model additionally incorporated Selective Distribution
Consistency (SDC) and Latent Consistency Matching (LCM). Both settings used
the same MTP architecture, allowing the effects of the consistency objectives
to be evaluated independently of architectural changes.

Unless otherwise specified, values reported with $\pm$ denote the
mean $\pm$ standard deviation over three random seeds.

The models were evaluated in terms of prediction quality, draft-token
acceptance, generation quality, and cross-dataset generalization.
The corresponding evaluation procedures and results are described
in the following sections.

\section{Main Experimental Results}\label{sec:4-main_results}

\subsection{Oracle Perplexity}

Oracle PPL was computed independently for each prediction head using
its corresponding ground-truth future token, where lower PPL indicates
better prediction quality.

\begin{table}[!htbp]
    \centering
    \caption{Oracle PPL comparison between the baseline and SDC+LCM.}
    \label{tab:oracle_ppl}
    \begin{tabular}{lcccc}
        \toprule[1.25pt]
        Setting & Head 1 & Head 2 & Head 3 & Head 4 \\
        \midrule
        Baseline
        & 19.4293 & 87.2351 & 200.0813 & 318.3247 \\

        SDC+LCM
        & \textbf{19.1695}
        & \textbf{83.7266}
        & \textbf{189.2825}
        & \textbf{299.5625} \\

        Improvement
        & 1.34\%
        & 4.02\%
        & 5.40\%
        & \textbf{5.89\%} \\
        \bottomrule[1.25pt]
    \end{tabular}
\end{table}

Table \ref{tab:oracle_ppl} compares the Oracle PPL of the baseline MTP model and the proposed SDC+LCM model across four prediction heads. The proposed method consistently reduces PPL for all prediction heads. Head 1 shows a relatively small improvement of 1.34\%, whereas larger improvements are observed for the later prediction heads. Specifically, the PPL reductions for Heads 2, 3, and 4 are 4.02\%, 5.40\%, and 5.89\%, respectively.

The improvement becomes more pronounced as the prediction position moves farther into the future, with Head 4 achieving the largest reduction. This result indicates that consistency-based training is particularly beneficial for the optimization of later prediction heads, which are responsible for more distant future-token predictions. Although the consistency objectives primarily constrain Heads 2–4, a slight improvement is also observed for Head 1.

\subsection{Draft Token Acceptance}\label{subsec:draft_token_acceptance}

\begin{table}[!htbp]
    \centering
    \caption{Comparison of Average Draft Tokens Accepted between the baseline and SDC+LCM.}
    \label{tab:draft_token_acceptance}
    \begin{tabular}{lccccc}
        \toprule[1.25pt]
        Setting
        & Overall
        & Prefix 128
        & Prefix 256
        & Prefix 512
        & Prefix 768 \\
        \midrule

        Baseline
        & 0.7010
        & 0.9117 $\pm$ 0.0205
        & 0.9458 $\pm$ 0.0068
        & 0.9689 $\pm$ 0.0221
        & 0.9760 $\pm$ 0.0222 \\

        SDC+LCM
        & \textbf{0.7302}
        & \textbf{0.9653 $\pm$ 0.0401}
        & \textbf{0.9873 $\pm$ 0.0063}
        & \textbf{1.0158 $\pm$ 0.0309}
        & \textbf{1.0272 $\pm$ 0.0250} \\

        Improvement
        & 4.17\%
        & \textbf{5.88\%}
        & 4.39\%
        & 4.84\%
        & 5.25\% \\
        \bottomrule[1.25pt]
    \end{tabular}
\end{table}

Draft-token acceptance was evaluated using greedy predictions from the
MTP heads. Head 1 was used as the verification reference, while predictions
from Heads 2--$K$ were treated as draft tokens. Draft tokens were accepted
sequentially when they matched the corresponding Head 1 predictions, and
verification stopped at the first mismatch. Head 1 itself was not included
in the accepted draft-token count.

Average Draft Tokens Accepted denotes the average number of consecutively
accepted draft tokens per evaluation instance. The overall result was computed
over all valid positions, while additional evaluations were conducted using
random contiguous prefixes of lengths 128, 256, 512, and 768.

Table \ref{tab:draft_token_acceptance} compares the Average Draft Tokens
Accepted between the baseline and SDC+LCM. The proposed method increases
the overall Average Draft Tokens Accepted from 0.7010 to 0.7302,
corresponding to an improvement of 4.17\%.

Consistent improvements are also observed across different prefix lengths.
At prefix lengths of 128, 256, 512, and 768, SDC+LCM achieves improvements
of 5.88\%, 4.39\%, 4.84\%, and 5.25\%, respectively. Among these settings,
the largest improvement is observed at a prefix length of 128, while the
highest absolute number of accepted draft tokens is achieved at a prefix
length of 768.

These results indicate that the proposed consistency objectives improve
the reliability of future-token predictions, resulting in more draft tokens
being accepted under the self-speculative decoding criterion.

\subsection{Generative Quality}\label{subsec:generative_quality}

Generation quality was evaluated using Generative PPL, entropy, and
4-gram repetition. A total of 200 samples were generated using the first
128 tokens of each evaluation sequence as the prompt. Generation was
performed using top-$p$ and temperature sampling with $p=0.9$ and a
temperature of 0.8. Generative PPL was computed using GPT-2 Large as an
external language model, where only the generated continuation was included
in the perplexity calculation. Entropy was used to characterize the
diversity of generated tokens, while 4-gram repetition was used to measure
repetitive generation patterns.

\begin{table}[!htbp]
    \centering
    \caption{Comparison of generative quality between the baseline and SDC+LCM.}
    \label{tab:generative_quality}
    \begin{tabular}{lccc}
        \toprule[1.25pt]
        Setting
        & Generative PPL
        & Entropy
        & 4-gram Repetition \\
        \midrule

        Baseline
        & 39.5580 $\pm$ 0.9425
        & 7.3223 $\pm$ 0.0351
        & 0.0229 $\pm$ 0.0084 \\

        SDC+LCM
        & \textbf{18.7180 $\pm$ 0.8276}
        & 6.8620 $\pm$ 0.0202
        & 0.0526 $\pm$ 0.0147 \\

        Change
        & \textbf{-52.68\%}
        & -6.29\%
        & +130.17\% \\
        \bottomrule[1.25pt]
    \end{tabular}
\end{table}

Table \ref{tab:generative_quality} compares the generation quality of the
baseline and SDC+LCM models. The proposed method substantially reduces
Generative PPL from 39.5580 to 18.7180, corresponding to a reduction of
52.68\%. This result indicates that the generated text is evaluated as
more predictable by the external language model.

In addition, the entropy decreases from 7.3223 to 6.8620, representing
a reduction of 6.29\%. However, the 4-gram repetition increases from
0.0229 to 0.0526, corresponding to a relative increase of 130.17\%.

These results reveal a trade-off in generation behavior. Although
SDC+LCM significantly improves Generative PPL, the lower entropy and
higher 4-gram repetition indicate reduced generation diversity and an
increased tendency toward repetitive patterns. Nevertheless, the
absolute increase in 4-gram repetition is approximately 0.03.

\section{Zero-Shot Evaluation}\label{sec:zero_shot_evaluation}

\begin{table}[!htbp]
    \centering
    \caption{Zero-shot Oracle PPL improvement of SDC+LCM over the baseline across different datasets.}
    \label{tab:zero_shot_ppl}
    \begin{tabular}{lcccc}
        \toprule[1.25pt]
        Dataset
        & Head 1
        & Head 2
        & Head 3
        & Head 4 \\
        \midrule

        WikiText
        & 0.32\%
        & \textbf{7.57\%}
        & 7.19\%
        & 5.62\% \\

        PTB
        & \textbf{13.82\%}
        & 9.90\%
        & 10.53\%
        & 5.63\% \\

        LAMBADA
        & \textbf{38.53\%}
        & 10.43\%
        & 4.44\%
        & 3.27\% \\

        arXiv
        & 1.48\%
        & \textbf{3.94\%}
        & 3.71\%
        & 0.27\% \\

        PubMed
        & 7.04\%
        & 7.76\%
        & \textbf{8.33\%}
        & 8.26\% \\

        \bottomrule[1.25pt]
    \end{tabular}
\end{table}

To evaluate the generalization ability of the proposed method across
different domains, zero-shot evaluation was conducted on WikiText\cite{merity2016wikitext},
Penn Treebank (PTB)\cite{marcus1993ptb}, LAMBADA\cite{paperno2016lambada}, and the arXiv and PubMed subsets of the Scientific Papers dataset\cite{cohan2018scientificpapers}.
Table \ref{tab:zero_shot_ppl} reports the relative Oracle PPL improvement
of SDC+LCM over the baseline for each prediction head.

Overall, SDC+LCM achieves PPL improvements across all evaluated datasets
and prediction heads, although the magnitude of improvement varies
considerably across domains. For Head 1, relatively large improvements
are observed on PTB\cite{marcus1993ptb}, LAMBADA\cite{paperno2016lambada}, and PubMed\cite{cohan2018scientificpapers}, with LAMBADA\cite{paperno2016lambada} showing the
largest reduction of 38.53\%. In contrast, the improvement on WikiText\cite{merity2016wikitext}
is limited to 0.32\%.

For WikiText\cite{merity2016wikitext}, the improvements are mainly concentrated on Heads 2--4,
with reductions of 7.57\%, 7.19\%, and 5.62\%, respectively, while the
improvement for Head 1 is relatively limited.

The improvement patterns also vary for later prediction heads across
different domains. For example, Head 4 achieves an 8.26\% improvement
on PubMed but only a 0.27\% improvement on arXiv. These results indicate
that the effectiveness of consistency-based training depends partly on
the characteristics of the evaluation domain, while generally yielding improvements across unseen datasets.

\section{Ablation Study}\label{sec:ablation_study}

\subsection{Overall Ablation Results}\label{subsec:overall_ablation_results}

\begin{table}[!htbp]
    \centering
    \caption{Ablation results of Oracle PPL under different consistency settings.}
    \label{tab:ablation_oracle_ppl}
    \small
    \begin{tabular}{lcccc}
        \toprule[1.25pt]
        Setting
        & Head 1
        & Head 2
        & Head 3
        & Head 4 \\
        \midrule

        Baseline
        & 26.2356 
        & 121.8033 
        & 269.2450 
        & 406.6959  \\

        $\lambda_{SDC}=0.5,\ \tau=0.3,\ \lambda_{LCM}=0$
        & 26.4587 
        & 120.7186 
        & 265.0660 
        & 402.0413  \\
        
        $\lambda_{SDC}=1.5,\ \tau=0.3,\ \lambda_{LCM}=0$
        & 27.2684 
        & 122.5244 
        & 264.6119 
        & 397.3106  \\

        \midrule

        $\lambda_{SDC}=1.0,\ \tau=0.0,\ \lambda_{LCM}=0$
        & 26.7541 
        & 120.6687 
        & 261.6664 
        & 394.6503  \\

        $\lambda_{SDC}=1.0,\ \tau=0.3,\ \lambda_{LCM}=0$
        & 26.8150 
        & 120.5276 
        & 261.9479 
        & 394.7055  \\

        $\lambda_{SDC}=1.0,\ \tau=0.6,\ \lambda_{LCM}=0$
        & 27.7396 
        & 129.2922
        & 278.8881 
        & 417.5874  \\

        \midrule

        $\lambda_{SDC}=0,\ \tau=0,\ \lambda_{LCM}=0.5$
        & 25.7096 
        & 118.9247 
        & 262.1637 
        & 397.9911  \\

        $\lambda_{SDC}=0,\ \tau=0,\ \lambda_{LCM}=1.0$
        & 25.5312 
        & 118.5823 
        & 261.1872 
        & 395.1838  \\

        $\lambda_{SDC}=0,\ \tau=0,\ \lambda_{LCM}=1.5$
        & \textbf{25.4583} 
        & 119.1595 
        & 263.2083 
        & 398.8679  \\

        \midrule

        $\lambda_{SDC}=1.0,\ \tau=0.3,\ \lambda_{LCM}=0.5$
        & 26.2402 
        & \textbf{118.2168} 
        & 257.0344 
        & 387.7455  \\

        $\lambda_{SDC}=1.0,\ \tau=0.3,\ \lambda_{LCM}=1.0$
        & 26.4553 
        & 119.2975 
        & 258.4100 
        & 390.7829  \\

        $\lambda_{SDC}=1.5,\ \tau=0.3,\ \lambda_{LCM}=1.0$
        & 26.5485
        & 118.4341 
        & \textbf{254.9612} 
        & \textbf{383.7626}  \\

        \bottomrule[1.25pt]
    \end{tabular}

    \vspace{2mm}
    
    \footnotesize{\textit{Bold values indicate the best result in each column.}}
\end{table}

\begin{table}[!htbp]
    \centering
    \caption{Ablation results of Average Draft Tokens Accepted under different consistency settings.}
    \label{tab:ablation_draft_tokens}
    \resizebox{\textwidth}{!}{
        \begin{tabular}{lccccc}
            \toprule[1.25pt]
            Setting
            & Overall
            & Prefix 128
            & Prefix 256
            & Prefix 512
            & Prefix 768 \\
            \midrule

            Baseline
            & 0.6238
            & 0.8251 $\pm$ 0.0140
            & 0.8205 $\pm$ 0.0062
            & 0.8519 $\pm$ 0.0195
            & 0.8552 $\pm$ 0.0233 \\

            $\lambda_{SDC}=0.5,\ \tau=0.3,\ \lambda_{LCM}=0$
            & 0.6345
            & 0.8195 $\pm$ 0.0122
            & 0.8416 $\pm$ 0.0115
            & 0.8601 $\pm$ 0.0065
            & 0.8644 $\pm$ 0.0304 \\

            $\lambda_{SDC}=1.5,\ \tau=0.3,\ \lambda_{LCM}=0$
            & 0.6414
            & 0.8522 $\pm$ 0.0098
            & 0.8581 $\pm$ 0.0139
            & 0.8897 $\pm$ 0.0043
            & \textbf{0.9082 $\pm$ 0.0204} \\

            \midrule

            $\lambda_{SDC}=1.0,\ \tau=0.0,\ \lambda_{LCM}=0$
            & 0.6388
            & 0.8374 $\pm$ 0.0165
            & 0.8746 $\pm$ 0.0108
            & 0.8736 $\pm$ 0.0158
            & 0.8802 $\pm$ 0.0029 \\

            $\lambda_{SDC}=1.0,\ \tau=0.3,\ \lambda_{LCM}=0$
            & 0.6400
            & 0.8406 $\pm$ 0.0131
            & 0.8602 $\pm$ 0.0232
            & 0.8720 $\pm$ 0.0233
            & 0.8772 $\pm$ 0.0015 \\

            $\lambda_{SDC}=1.0,\ \tau=0.6,\ \lambda_{LCM}=0$
            & 0.6113
            & 0.7908 $\pm$ 0.0074
            & 0.7919 $\pm$ 0.0086
            & 0.8009 $\pm$ 0.0035
            & 0.8174 $\pm$ 0.0130 \\

            \midrule

            $\lambda_{SDC}=0,\ \tau=0,\ \lambda_{LCM}=0.5$
            & 0.6398
            & 0.8517 $\pm$ 0.0156
            & 0.8589 $\pm$ 0.0096
            & 0.8794 $\pm$ 0.0050
            & 0.8895 $\pm$ 0.0257 \\

            $\lambda_{SDC}=0,\ \tau=0,\ \lambda_{LCM}=1.0$
            & 0.6442
            & 0.8491 $\pm$ 0.0130
            & 0.8698 $\pm$ 0.0154
            & 0.8919 $\pm$ 0.0168
            & 0.8943 $\pm$ 0.0012 \\

            $\lambda_{SDC}=0,\ \tau=0,\ \lambda_{LCM}=1.5$
            & 0.6467
            & 0.8510 $\pm$ 0.0112
            & \textbf{0.8862 $\pm$ 0.0054}
            & 0.8951 $\pm$ 0.0129
            & 0.9028 $\pm$ 0.0327 \\

            \midrule

            $\lambda_{SDC}=1.0,\ \tau=0.3,\ \lambda_{LCM}=0.5$
            & 0.6507
            & 0.8498 $\pm$ 0.0339
            & 0.8615 $\pm$ 0.0179
            & 0.8987 $\pm$ 0.0103
            & 0.8862 $\pm$ 0.0151 \\

            $\lambda_{SDC}=1.0,\ \tau=0.3,\ \lambda_{LCM}=1.0$
            & 0.6495
            & 0.8402 $\pm$ 0.0209
            & 0.8612 $\pm$ 0.0043
            & 0.8879 $\pm$ 0.0157
            & 0.8950 $\pm$ 0.0071 \\

            $\lambda_{SDC}=1.5,\ \tau=0.3,\ \lambda_{LCM}=1.0$
            & \textbf{0.6525}
            & \textbf{0.8590 $\pm$ 0.0101}
            & 0.8656 $\pm$ 0.0214
            & \textbf{0.9023 $\pm$ 0.0277}
            & 0.9044 $\pm$ 0.0140 \\

            \bottomrule[1.25pt]
        \end{tabular}
    }
    \vspace{2mm}
    
    \footnotesize{\textit{Bold values indicate the best result in each column.}}
\end{table}

\begin{table}[!htbp]
    \centering
    \caption{Ablation results of generative quality under different consistency settings.}
    \label{tab:ablation_generation}
    \small
    \begin{tabular}{lccc}
        \toprule[1.25pt]
        Setting
        & Generative PPL
        & Entropy
        & 4-gram Repetition \\
        \midrule

        Baseline
        & 83.8285 $\pm$ 5.1409
        & 7.4068 $\pm$ 0.0461
        & 0.0119 $\pm$ 0.0097 \\

        $\lambda_{SDC}=0.5,\ \tau=0.3,\ \lambda_{LCM}=0$
        & 52.7654 $\pm$ 2.2732
        & 7.1490 $\pm$ 0.0335
        & 0.0323 $\pm$ 0.0169 \\

        $\lambda_{SDC}=1.5,\ \tau=0.3,\ \lambda_{LCM}=0$
        & 31.8082 $\pm$ 0.4503
        & 6.8248 $\pm$ 0.0353
        & 0.0485 $\pm$ 0.0159 \\

        \midrule

        $\lambda_{SDC}=1.0,\ \tau=0.0,\ \lambda_{LCM}=0$
        & 68.0596 $\pm$ 4.5760
        & 7.3141 $\pm$ 0.0519
        & 0.0193 $\pm$ 0.0154 \\

        $\lambda_{SDC}=1.0,\ \tau=0.3,\ \lambda_{LCM}=0$
        & 48.4875 $\pm$ 1.2280
        & 7.1184 $\pm$ 0.0336
        & 0.0183 $\pm$ 0.0055 \\

        $\lambda_{SDC}=1.0,\ \tau=0.6,\ \lambda_{LCM}=0$
        & 28.2332 $\pm$ 1.2754
        & 6.7789 $\pm$ 0.0178
        & 0.0393 $\pm$ 0.0080 \\

        \midrule

        $\lambda_{SDC}=0,\ \tau=0,\ \lambda_{LCM}=0.5$
        & 26.2091 $\pm$ 1.5675
        & 6.8576 $\pm$ 0.0374
        & 0.0555 $\pm$ 0.0204 \\

        $\lambda_{SDC}=0,\ \tau=0,\ \lambda_{LCM}=1.0$
        & 24.0955 $\pm$ 3.2440
        & 6.7935 $\pm$ 0.0828
        & 0.0657 $\pm$ 0.0388 \\

        $\lambda_{SDC}=0,\ \tau=0,\ \lambda_{LCM}=1.5$
        & 29.9033 $\pm$ 0.5373
        & 6.8644 $\pm$ 0.0249
        & 0.0215 $\pm$ 0.0049 \\

        \midrule

        $\lambda_{SDC}=1.0,\ \tau=0.3,\ \lambda_{LCM}=0.5$
        & \textbf{21.4788 $\pm$ 0.0773}
        & 6.6876 $\pm$ 0.0421
        & 0.0720 $\pm$ 0.0062 \\

        $\lambda_{SDC}=1.0,\ \tau=0.3,\ \lambda_{LCM}=1.0$
        & 26.3414 $\pm$ 0.3656
        & 6.7567 $\pm$ 0.0016
        & 0.0460 $\pm$ 0.0160 \\

        $\lambda_{SDC}=1.5,\ \tau=0.3,\ \lambda_{LCM}=1.0$
        & 25.7832 $\pm$ 1.6338
        & 6.7177 $\pm$ 0.0455
        & 0.0453 $\pm$ 0.0197 \\

        \bottomrule[1.25pt]
    \end{tabular}

    \vspace{2mm}
    
    \footnotesize{\textit{Bold indicates the lowest Generative PPL.}}
\end{table}

\FloatBarrier

\subsection{Effect of SDC Weight}\label{subsece:effect_of_sdc_weight}

To investigate the effect of the SDC loss weight, the confidence threshold
and LCM weight were fixed at $\tau=0.3$ and $\lambda_{LCM}=0$, respectively, while
the SDC weight was varied among $\lambda_{SDC}=0.5$, $1.0$, and $1.5$. The baseline
without consistency regularization is also included for comparison.

Increasing the SDC weight generally improves draft-token acceptance and
reduces Generative PPL, while the Oracle PPL results indicate that excessively
strong distribution consistency does not consistently benefit all prediction
heads.

\subsection{Effect of Confidence Threshold}\label{subsece:effect_of_confidence_threshold}

To investigate the effect of the confidence threshold, $\lambda_{SDC}=1.0$ and
$\lambda_{LCM}=0$ were fixed while $\tau$ was varied among 0.0, 0.3, and 0.6.

The Oracle PPL results are similar between $\tau=0.0$ and $\tau=0.3$,
whereas $\tau=0.6$ significantly degrades all prediction heads. A similar
trend is observed for Average Draft Tokens Accepted, which decreases from
0.6400 at $\tau=0.3$ to 0.6113 at $\tau=0.6$.

In contrast, Generative PPL continues to decrease as $\tau$ increases,
but this is accompanied by lower entropy and higher 4-gram repetition.
Therefore, an excessively high confidence threshold may improve
Generative PPL while degrading future-token prediction and draft-token
acceptance, indicating a trade-off among evaluation metrics.

\subsection{Effect of LCM Weight}\label{subsece:effect_of_lcm_weight}

To investigate the effect of the LCM weight, $\lambda_{SDC}=0$ was fixed while
$\lambda_{LCM}$ was varied among 0.5, 1.0, and 1.5.

Increasing $\lambda_{LCM}$ from 0.5 to 1.0 generally improves Oracle PPL across
the prediction heads. However, further increasing $\lambda_{LCM}$ to 1.5 slightly
degrades the PPL of Heads 2--4, suggesting that excessive latent
consistency does not provide additional benefits.

Average Draft Tokens Accepted gradually increases from 0.6398 to 0.6467
as $\lambda_{LCM}$ increases. In contrast, Generative PPL reaches its lowest value
at $\lambda_{LCM}=1.0$ and increases again at $\lambda_{LCM}=1.5$. 
These results indicate that a moderate LCM weight provides a better
balance between prediction quality and draft-token acceptance.

\subsection{Joint Effect of SDC and LCM}\label{subsece:effect_of_joint_effect}

To evaluate the joint effect of SDC and LCM, configurations using both
consistency objectives were compared with SDC-only and LCM-only settings.

Joint training generally achieves lower Oracle PPL for Heads 2--4 and
higher Average Draft Tokens Accepted than using either consistency
objective alone. In particular, the joint configurations achieve draft
acceptance values of approximately 0.65, compared with 0.6400 for the
SDC-only setting and 0.6442 for the LCM-only setting.

For generation quality, joint training further reduces Generative PPL,
but is accompanied by lower entropy and higher 4-gram repetition.
Therefore, SDC and LCM can complement each other in improving future-token
prediction and draft-token acceptance, while a trade-off with generation
diversity remains.

\section{Discussion}\label{subsece:discussion}

The experimental results show that consistency-based training is particularly
beneficial for prediction heads responsible for more distant future tokens.
Across the main experiments, SDC+LCM consistently reduces Oracle PPL, with
larger improvements generally observed for the later prediction heads.

The improvement in future-token prediction is also reflected in draft-token
acceptance. SDC+LCM increases the Average Draft Tokens Accepted, while the
zero-shot results indicate that the benefits of consistency-based training
can generalize across different domains.

However, the ablation study shows that stronger consistency constraints do
not always lead to better overall performance. In particular, lower
Generative PPL is often accompanied by lower entropy and higher repetition,
indicating a trade-off between prediction consistency and generation
diversity. Therefore, the strengths of SDC, LCM, and the confidence threshold
should be carefully balanced.
